\documentclass[11pt]{article}
\input{_report_preamble}

\title{\textbf{Cycloidal Gearbox: Dynamometer Test Report}}
\author{Paladin Engineering}
\date{\today}

\begin{document}
\maketitle

\resultfig{pics/gearboxCropped.png}{0.4}{Cycloidal-drive gearbox design that was tested (Note: we were provided 2 gearboxes.  This document concerns gearbox 1, and this picture is of gearbox 2.)}{fig:gearbox}

% ===========================================================================
\section{Introduction}
% ===========================================================================
This report summarizes dynamometer characterization of the cycloidal reduction
gearbox driven by a test motor. The unit is a \textbf{26:1}
cycloidal drive, tested input-driven and (where noted) output-driven on the
Paladin dynamometer with in-line torque cells on both the input and output shafts.

A cycloidal drive transmits torque through a lobed cycloidal disc rolling inside
a ring of pins, with multiple lobes engaged at once. That architecture drives the
properties we measure:

\begin{itemize}
  \item \textbf{Low backlash, high stiffness.} Many simultaneous, preloaded
        rolling contacts mean very little lost motion and a stiff torsional path
        (Section~\ref{sec:backlash}).
  \item \textbf{High, load-dependent efficiency.} Rolling (rather than sliding)
        contact gives good efficiency that improves with load as fixed losses
        become a smaller fraction of throughput (Section~\ref{sec:eff}).
  \item \textbf{Low no-load drag with lobe-frequency ripple.} The disc's orbiting
        motion produces a small torque ripple at the lobe-pass frequency as rolling contacts hand off, riding
        on top of a low baseline drag (Section~\ref{sec:gbrunning}).
  \item \textbf{Good backdrivability.} The same low friction makes the unit
        readily back-driven, at a modest efficiency penalty vs.\ forward driving
        (Section~\ref{sec:eff}).
\end{itemize}

Tests reported: backlash/torsional analysis, an efficiency map (forward and
backdrive), no-load running torque, and stiction. A structural tap test is
included as an appendix where it bears on the dynamic results.

% ===========================================================================
\section{Backlash}\label{sec:backlash}
% ===========================================================================
The output side of the shaft was held at a fixed position while a torque sawtooth was applied at the
input; the output was then shifted and the sweep repeated, sampling backlash at
19 mesh positions. Referring the input angle to the output side
($\theta_\text{out}-\theta_\text{in}/26$) and plotting deflection vs.\ tared input
torque yields a hysteresis loop per position. Very little apparent backlash is present,
with results looking more similar to a stiff torsional spring with ripple effects than a traditional dead-band.
It is likely that this mechanism is preloaded, rendering a traditional flank-based backlash measurement less meaningful.

\resultfig{gearbox_tests/results/backlash/backlash_deflection_vs_torque.png}{0.78}
  {Output-referred deflection vs.\ tared input torque across 19 mesh positions
   (color = output hold position).}{fig:backlash}

The loop is narrow and the loaded flanks are nearly straight and almost
collinear, which is the clean signature of a \textbf{preloaded drive with minimal
backlash}: the few arc-minutes of gap are small and consistent around the output
rotation, and the high flank slope confirms a stiff torsional path. The
position-to-position scatter (color bands) reflects mesh-phase variation rather
than any single large dead-band.

Torsional stiffness is high, with a mean of \SI{10580}{\newton\meter\per\radian} and
  a standard deviation of \SI{540}{\newton\meter\per\radian} across the 19 mesh positions. The stiffness is not constant, but varies with mesh position and torque, as expected for a cycloidal drive with multiple rolling contacts.

% ===========================================================================
\section{Efficiency}\label{sec:eff}
% ===========================================================================
The input was velocity-controlled while the load applied an output torque
setpoint; both shaft torques were measured at each $\sim\!\SI{1}{\second}$ dwell.
Because the 26:1 ratio is kinematically fixed (no slip), power efficiency equals
torque efficiency: $\eta = \min(26|T_\text{in}|,|T_\text{out}|)/\max(26|T_\text{in}|,|T_\text{out}|)$.
The sign of $\omega_\text{in}\!\cdot\!T_\text{out}$ splits the data into forward
driving and backdriving maps; low-torque dwells (where $\eta$ is noise) are
excluded.

\begin{figure}[H]
  \centering
  \includegraphics[width=0.49\linewidth]{gearbox_tests/results/efficiency/efficiency_map_forward.png}
  \includegraphics[width=0.49\linewidth]{gearbox_tests/results/efficiency/efficiency_map_backdrive.png}
  \caption{Efficiency maps vs.\ output speed and output torque: forward driving
  (left) and backdriving (right), CW/CCW averaged.}
  \label{fig:eff}
\end{figure}

\begin{table}[H]
  \centering
  \caption{Efficiency summary (loaded operating points).}
  \begin{tabular}{lrr}
    \toprule
    & Forward & Backdrive\\
    \midrule
    Peak efficiency  & \SI{95.4}{\percent} & \SI{93.4}{\percent}\\
    Mean (loaded)    & \SI{93.1}{\percent} & \SI{90.1}{\percent}\\
    \bottomrule
  \end{tabular}
\end{table}

Efficiency \textbf{rises with output torque}, from the high 70s/80s\,\% at
light load up to the low-to-mid 90s\,\%, because the roughly fixed friction
loss is a shrinking fraction of the transmitted torque. It falls off slightly at
the highest input speeds. Backdriving sits a few points below forward driving at
matched operating points, consistent with the overhauling direction paying the
same friction loss against a smaller delivered torque. These current plots are fairly speed-limited, and a future test campaign will explore the full speed range of the gearbox.

% ===========================================================================
\section{Stiction}\label{sec:stiction}
% ===========================================================================
Static-friction (breakaway) torque was measured with a slow torque-ramp
protocol: the test motor was commanded a gentle triangular torque sweep
($\sim\!\SI{0.2}{\newton\meter\per\second}$) that rises until the shaft breaks
free and slips, then reverses; the torque held just before first motion is the
breakaway torque. Two configurations were run: the
gearbox driven from its \textbf{input} shaft, and (with the gearbox flipped in
place) driven from its \textbf{output} shaft.

The reported gearbox stiction is read from the in-line torque cell.
Because the cell sits at the gearbox shaft, it measures the torque actually
delivered into the gearbox at breakaway. This is the gearbox's own static friction,
already free of the motor's internal friction, with no subtraction required to account for drive motor friction.

\resultfig{gearbox_tests/results/stiction/stiction_comparison.png}{0.7}
  {Breakaway (static-friction) torque measured by the in-line torque cell at the
   gearbox shaft: input-driven (at the input shaft) vs.\ output-driven (at the
   output shaft).}{fig:stiction}

Input-shaft stiction is very small ($\approx\SI{0.04}{\newton\meter}$ at the
input), consistent with the lightly-loaded, predominantly-rolling high-speed input
stage. The output-driven breakaway is far larger and much noisier; its large
spread likely reflects cycloidal lobe-pass ripple rather than poor repeatability.

% ===========================================================================
\section{Running Torque}\label{sec:gbrunning}
% ===========================================================================
With the opposite shaft free, the driven shaft was swept through a $\pm$velocity
triangle and the no-load drag (Coulomb~$+$~viscous) extracted by velocity-binning
to cancel inertia. Both data sets are plotted against \emph{output} speed: the
input-driven sweep is referred to the output through the fixed 26:1 ratio
($\omega_\text{out}=\omega_\text{in}/26$), whereas the output-driven sweep (gearbox
flipped) already measures the output shaft directly and is \emph{not} rescaled.

\subsection{Input-driven}
\resultfig{gearbox_tests/results/running_torque/running_torque_drag_vs_velocity.png}{0.78}
  {No-load input running torque vs.\ output speed, with the Coulomb+viscous fit.}{fig:gbrun-in}

\begin{table}[H]
  \centering
  \caption{Input-driven no-load drag.}
  \begin{tabular}{lr}
    \toprule
    Quantity & Value\\
    \midrule
    Coulomb drag         & \SI{0.040}{\newton\meter}\\
    Viscous drag         & \SI{0.0328}{\newton\meter\per(\radian\per\second)}\\
    Drag at max speed (\SI{2.4}{\radian\per\second}) & \SI{0.120}{\newton\meter}\\
    \bottomrule
  \end{tabular}
\end{table}

Input-referred drag is very low and
grows only weakly with speed, again indicating predominantly rolling contact. A
torque ripple rides on top of this baseline at the cycloidal lobe-pass frequency
and grows large at the highest input speeds, but that growth is mostly the test
fixture's structural resonance being excited rather than the gearbox itself (see
Figure~\ref{fig:gbrun-fft}), so its raw amplitude here is not a reliable ripple
figure. The ripple magnitude is instead quantified under load below
(Figure~\ref{fig:ripple}).

\subsection{Output-driven}
With the gearbox flipped in place, the test motor drove the output shaft directly,
so the measured speed is already output speed and the in-line torque is the
output-referred back-drive drag.

\resultfig{gearbox_tests/results/output_running_torque/output_running_torque_drag_vs_velocity.png}{0.78}
  {No-load output-driven (back-drive) running torque vs.\ output speed, with the
   Coulomb+viscous fit.}{fig:gbrun-out}

\begin{table}[H]
  \centering
  \caption{Output-driven (back-drive) no-load drag, output-referred.}
  \begin{tabular}{lr}
    \toprule
    Quantity & Value\\
    \midrule
    Coulomb drag         & \SI{1.51}{\newton\meter}\\
    Viscous drag         & \SI{0.191}{\newton\meter\per(\radian\per\second)}\\
    Drag at max speed (\SI{10.3}{\radian\per\second}) & \SI{3.48}{\newton\meter}\\
    \bottomrule
  \end{tabular}
\end{table}

Driven directly on the output, the motor's speed range maps straight to output
speed rather than being divided by 26, so this sweep reaches far higher output
speed (\SI{10.3}{\radian\per\second} vs.\ \SI{2.4}{}) and exercises the viscous
term much more strongly. The two runs
are mutually consistent: referring the input-driven Coulomb drag to the output by
the power-equivalent factor of 26 ($0.040\times26\approx\SI{1.0}{\newton\meter}$)
lands close to the directly-measured back-drive value of \SI{1.51}{\newton\meter},
the modest excess being the expected back-driving friction penalty.

\subsection{Frequency content and torque ripple}
\resultfig{gearbox_tests/results/output_running_torque/output_running_torque_fft.png}{0.85}
  {Output-driven running-torque spectrum (top, segment-averaged) and spectrogram
   (bottom). The input-driven content is similar and is not shown separately.}{fig:gbrun-fft}

This section is the frequency-domain view of the \textbf{torque ripple}. The drag
curves above describe only the mean (baseline) torque; subtracting that baseline
leaves the small periodic fluctuation that rides on top of it, and the spectrum and
spectrogram below sort that fluctuation by frequency to show where the ripple
energy actually sits. The ripple grows with speed because the same per-revolution
disturbance is swept to higher frequency (and amplified near resonances) as the
shaft turns faster.

The spectrogram resolves the ripple into speed-proportional \emph{orders}, the
diagonal ridges that sweep up in frequency as the shaft accelerates, which are
the cycloidal lobe-pass forcing and its harmonics. Riding on top of these is a
broad, \emph{speed-independent} peak near \SIrange{18}{24}{\hertz} in the averaged
spectrum. Because it stays at a fixed frequency while the speed ramps, it is a
\textbf{structural resonance} of the test fixture, not a gearbox order: it
coincides with the \SI{21.4}{\hertz} fundamental from the tap test
(Appendix~\ref{sec:tap}), with the lobe-pass forcing ringing the fixture as it
sweeps through.

A clean ripple \emph{magnitude} is best taken not from these no-load sweeps but
from the loaded, constant-speed efficiency dwells (Section~\ref{sec:eff}): with the
output torque held at a fixed setpoint, whatever fluctuation remains on it is the
transmission ripple, measured at a real load and below the fixture resonance.

\resultfig{gearbox_tests/results/torque_ripple/torque_ripple.png}{0.72}
  {Per-dwell output torque ripple (RMS) vs.\ transmitted torque. Rotating dwells
   (blue) sit well above the no-rotation static-hold floor (grey).}{fig:ripple}

The output torque ripples by \textbf{$\approx\SI{0.9}{\newton\meter}$ RMS}, about
\SI{1.3}{\percent} of the transmitted torque. Two features mark it as a genuine
gearbox effect: it is nearly the \emph{same} absolute size at 50, 80, and
\SI{110}{\newton\meter} of load (so it reads as a larger percentage at light load),
and it drops to a $\approx\SI{0.3}{\newton\meter}$ floor when the shaft is held
still, showing it is produced by rotation (the cycloidal lobe-pass hand-off)
rather than measurement noise.

% ===========================================================================
\appendix
\section{Structural Tap Test}\label{sec:tap}
% ===========================================================================
The stationary, powered dynamometer was struck with a rubber mallet and the
input-torque ring-down was FFT'd to find structural natural frequencies.

\resultfig{gearbox_tests/results/tap_test/tap_test_fft.png}{0.78}
  {Impulse ring-down spectrum of the input-torque signal.}{fig:tap}

The dominant mode is at \SI{21.4}{\hertz}, with further peaks at \SI{64}{\hertz}
and a cluster near \SI{407}{}--\SI{438}{\hertz}. This appendix is included because
the \SI{21.4}{\hertz} mode lines up with the broad structural peak in the
running-torque spectrum (Section~\ref{sec:gbrunning}): that peak is the fixture
resonating as the cycloidal lobe-pass forcing sweeps through its natural frequency,
not a property of the gearbox itself. It is worth keeping in mind when interpreting
the ripple seen in the running-torque sweeps.

\end{document}
